\documentclass[11pt]{article}
\usepackage[margin=1in]{geometry}
\usepackage{amsmath,amsthm,amssymb}
\usepackage[hidelinks]{hyperref}
\newtheorem{theorem}{Theorem}[section]
\newtheorem{lemma}[theorem]{Lemma}
\newtheorem{corollary}[theorem]{Corollary}
\theoremstyle{remark}
\newtheorem{remark}[theorem]{Remark}
\newtheorem{example}[theorem]{Example}
\newcommand{\Int}{\operatorname{Int}}
\newcommand{\F}{\mathbb F}
\newcommand{\N}{\mathbb N}
\newcommand{\Nzero}{\mathbb N_0}
\title{Divisors and unique factorizations of Carlitz binomial polynomials}
\author{}
\date{September 8, 2026}

\begin{document}
\maketitle
\begin{abstract}
Let $A=\F_q[t]$. We classify the divisors in $\Int(A)$ of every finite product of the basic Carlitz binomial polynomials $\beta_{q^s}$. Each divisor is, up to a unit, a subproduct of the given basic factors. Thus the multiplicative monoid generated by these factors is divisor-closed and factorial. In particular, every power of every Carlitz binomial polynomial has a unique factorization into irreducibles, answering a question of Tichy and Windisch. The proof combines a simultaneous-minimum identity for weighted products over the sets of polynomials of bounded degree with an extraction argument at a prime of prescribed degree.
\end{abstract}

\section{Introduction and main theorem}

Fix a prime power $q$, and write $A=\F_q[t]$ and $K=\F_q(t)$. The ring of integer-valued polynomials on $A$ is
\[
 \Int(A)=\{H\in K[X]:H(A)\subseteq A\}.
\]
Here $t$ and $X$ are independent indeterminates. Carlitz constructed a regular $A$-module basis of this ring that is analogous to the ordinary binomial basis of $\Int(\mathbb Z)$; we follow the notation and normalization in Tichy--Windisch~\cite{TW}.

Set $V_0=\{0\}$, $\psi_0(X)=X$, and $F_0=1$. For $s\geq1$, let
\begin{equation}\label{eq:def}
 \begin{split}
 V_s&=\{f\in A:\deg f<s\},\qquad
 \psi_s(X)=\prod_{f\in V_s}(X-f),\\
 F_s&=\prod_{i=1}^s(t^{q^i}-t)^{q^{s-i}},\qquad
 B_s(X)=\frac{\psi_s(X)}{F_s}.
 \end{split}
\end{equation}
We also put $B_0=X$. The convention $\deg0=-\infty$ is used. If $k=\sum_{i\geq0}\alpha_iq^i$ is the base-$q$ expansion of a nonnegative integer, then its Carlitz binomial polynomial is
\begin{equation}\label{eq:beta}
 \beta_k(X)=\prod_{i\geq0}B_i(X)^{\alpha_i}.
\end{equation}
In particular, $\beta_{q^s}=B_s$.

Tichy and Windisch proved that the $B_s$ are irreducible in $\Int(A)$ and that each is absolutely irreducible: every power of an individual $B_s$ factors uniquely into irreducibles~\cite[Theorem~1.2]{TW}. Their Question~3.1 asks whether this uniqueness persists for powers of the general polynomial~\eqref{eq:beta}. We prove the following stronger divisor statement.

\begin{theorem}\label{thm:main}
Let $a_0,\ldots,a_s\in\Nzero$, and put
\[
 P=\prod_{i=0}^s B_i^{a_i}.
\]
An element $H\in\Int(A)$ divides $P$ in $\Int(A)$ if and only if
\begin{equation}\label{eq:divisor}
 H=u\prod_{i=0}^sB_i^{b_i}
 \quad\text{for some }u\in\F_q^\times
 \text{ and integers }0\leq b_i\leq a_i.
\end{equation}
The unit $u$ and the exponent vector $(b_i)$ are unique.
\end{theorem}

\begin{corollary}\label{cor:question}
For every $k\geq1$ and $m\geq1$, the polynomial $\beta_k^m$ has a unique factorization into irreducibles in $\Int(A)$, up to order and associates. If $k=\sum_i\alpha_iq^i$, this factorization is
\[
 \beta_k^m=\prod_iB_i^{m\alpha_i}.
\]
\end{corollary}

The mixed-product assertion is stronger than absolute irreducibility of the separate $B_i$. The proof must show that extracting a highest-level basic factor from an arbitrary divisor preserves integer-valuedness at every prime. Section~2 provides the required fixed-divisor identity; Section~3 performs the extraction and proves the theorem. The argument also recovers the known irreducibility of the $B_i$.

\section{A simultaneous-minimum identity}

For a monic irreducible polynomial $\pi\in A$, let $v_\pi$ be the corresponding valuation on $K$, normalized by $v_\pi(\pi)=1$, with $v_\pi(0)=+\infty$. If $R\in A[X]$ is nonzero, write
\[
 \mu_\pi(R)=\min_{x\in A}v_\pi(R(x)).
\]
This minimum is finite: $A$ is infinite, and $R$ has only finitely many roots. It is the $\pi$-adic valuation of the fixed divisor of $R$.

The residue counts for $V_s$ and the formula for $v_\pi(F_s)$ used below are classical; see Wagner~\cite[p.~51]{Wagner}. We require a point that simultaneously minimizes the valuation of $\psi_s$ and of an arbitrary product of its linear factors with nonnegative multiplicities.

\begin{lemma}\label{lem:min}
Let $s\geq1$, let $\pi\in A$ be monic irreducible of degree $d$, and put
\[
 E_s(\pi)=\sum_{j=1}^{\lfloor s/d\rfloor}q^{s-jd}.
\]
Let
\[
 R(X)=\prod_{f\in V_s}(X-f)^{r_f},\qquad r_f\in\Nzero.
\]
Then, for every $c\in\Nzero$,
\begin{equation}\label{eq:min}
 \mu_\pi(\psi_s^cR)=cE_s(\pi)+\mu_\pi(R).
\end{equation}
Moreover,
\begin{equation}\label{eq:fixed}
 \mu_\pi(\psi_s)=E_s(\pi)=v_\pi(F_s).
\end{equation}
\end{lemma}

\begin{proof}
Write $s=kd+r$, where $0\leq r<d$. For $1\leq j\leq k$, the reduction map
\[
 V_s\longrightarrow A/(\pi^j)
\]
is surjective, and each fiber has $q^{s-jd}$ elements. Indeed, division by the monic polynomial $\pi^j$ expresses a polynomial of degree less than $s$ uniquely as a remainder of degree less than $jd$ plus $\pi^j$ times a polynomial of degree less than $s-jd$. Consequently, for every $x\in A$,
\begin{equation}\label{eq:trunc}
 \begin{split}
 \sum_{f\in V_s}\min\{v_\pi(x-f),k\}
 &=\sum_{j=1}^k\#\{f\in V_s:x\equiv f\pmod{\pi^j}\}\\
 &=E_s(\pi).
 \end{split}
\end{equation}
Empty sums have value zero, so this also holds for $k=0$.

The residue class $x+\pi^kA$ contains exactly $q^r$ elements of $V_s$. For $k=0$ this says simply that $|V_s|=q^s=q^r$. These elements occupy at most $q^r$ of the $q^d$ residue classes modulo $\pi^{k+1}$ lying above $x$ modulo $\pi^k$. Since $q^r<q^d$, choose $y\in A$ in one of the unoccupied classes. Then $y\equiv x\pmod{\pi^k}$, and for every $f\in V_s$,
\begin{equation}\label{eq:simultaneous}
 v_\pi(y-f)=\min\{v_\pi(x-f),k\}.
\end{equation}
For valuations less than $k$, this follows from the congruence between $x$ and $y$; for valuations at least $k$, it follows from the choice of the class of $y$ modulo $\pi^{k+1}$.

Equation~\eqref{eq:trunc} now gives
\[
 v_\pi(\psi_s(y))=E_s(\pi),
 \qquad v_\pi(R(y))\leq v_\pi(R(x)).
\]
It also shows that $v_\pi(\psi_s(x))\geq E_s(\pi)$ for all $x\in A$. Start with an $x$ attaining $\mu_\pi(R)$. The resulting $y$ attains this same minimum and also minimizes $\psi_s$. Hence
\[
 \mu_\pi(\psi_s^cR)=cE_s(\pi)+\mu_\pi(R).
\]
Taking $R=1$ and $c=1$ proves the first equality in~\eqref{eq:fixed}.

The polynomial $\pi$ divides $t^{q^i}-t$ exactly when $d\mid i$, and then its multiplicity is one. This follows from the description of the roots of $t^{q^i}-t$ as $\F_{q^i}$ and from its derivative $-1$. Applying $v_\pi$ to the product defining $F_s$ gives the second equality in~\eqref{eq:fixed}.
\end{proof}

\begin{corollary}\label{cor:criterion}
Every $B_s$ belongs to $\Int(A)$. More generally, if $C\in K^\times$, $s\geq1$, $c\in\Nzero$, and $R$ is as in Lemma~\ref{lem:min}, then
\begin{equation}\label{eq:criterion}
 C\psi_s^cR\in\Int(A)
 \quad\Longleftrightarrow\quad
 CF_s^cR\in\Int(A).
\end{equation}
\end{corollary}

\begin{proof}
An element $z\in K$ belongs to $A$ if and only if $v_\pi(z)\geq0$ for every monic irreducible $\pi\in A$. Therefore, for a nonzero polynomial $Q\in A[X]$,
\[
 CQ\in\Int(A)
 \quad\Longleftrightarrow\quad
 v_\pi(C)+\mu_\pi(Q)\geq0\quad\text{for every }\pi.
\]
Lemma~\ref{lem:min} makes these inequalities identical for the two polynomials in~\eqref{eq:criterion}. It also proves $\psi_s/F_s\in\Int(A)$. The case $B_0=X$ is immediate. All primes in the criterion are finite primes of $A$; no condition at the infinite place is imposed.
\end{proof}

\section{Extraction and divisor classification}

\begin{lemma}\label{lem:extract}
Let $s\geq1$, let $a_0,\ldots,a_s\in\Nzero$ with $a_s>0$, and suppose $H,J\in\Int(A)$ satisfy
\[
 HJ=\prod_{i=0}^sB_i^{a_i}.
\]
There is an integer $c$ with $0\leq c\leq a_s$ such that
\[
 H/B_s^c\in\Int(A),\qquad
 J/B_s^{a_s-c}\in\Int(A).
\]
Both quotients split over $K$ with all roots in $V_{s-1}$.
\end{lemma}

\begin{proof}
Unique factorization in $K[X]$ gives representations
\[
 H=C\prod_{f\in V_s}(X-f)^{u_f},\qquad
 J=D\prod_{f\in V_s}(X-f)^{w_f-u_f},
\]
where $C,D\in K^\times$, $0\leq u_f\leq w_f$, and $w_f$ is the multiplicity of $f$ in the numerator of $\prod_iB_i^{a_i}$. In particular,
\begin{equation}\label{eq:outer}
 w_f=a_s\qquad(f\in V_s\setminus V_{s-1}).
\end{equation}

Choose a monic irreducible polynomial $\pi\in A$ of degree $s$. Such a polynomial exists over every finite field in every positive degree. The set $V_s$ is a complete system of representatives modulo $\pi$, and Lemma~\ref{lem:min} gives
\[
 v_\pi(F_s)=1,\qquad v_\pi(F_i)=0\quad(0\leq i<s).
\]
Set $c=-v_\pi(C)$. Comparing leading coefficients in the identity for $HJ$ gives
\[
 -v_\pi(D)=a_s-c.
\]
For each $f\in V_s$, evaluate at $x=f+\pi$. Then $v_\pi(x-f)=1$, while $v_\pi(x-g)=0$ for every $g\in V_s\setminus\{f\}$. The integer-valuedness of $H$ and $J$ yields
\begin{equation}\label{eq:bounds}
 u_f\geq c,\qquad w_f-u_f\geq a_s-c.
\end{equation}
For $f$ in the nonempty set $V_s\setminus V_{s-1}$, equations~\eqref{eq:outer} and~\eqref{eq:bounds} force $u_f=c$. Since $0\leq u_f\leq a_s$ on this set, we have $0\leq c\leq a_s$.

For all $f\in V_s$, the multiplicities $u_f-c$ and $w_f-u_f-(a_s-c)$ are nonnegative. They are zero on $V_s\setminus V_{s-1}$. We can consequently write
\[
 H=C\psi_s^cR,
\]
where $R$ has nonnegative root multiplicities and all roots in $V_{s-1}$. Regard $R$ as a weighted product over $V_s$ with zero weights on the outer set. Corollary~\ref{cor:criterion} gives
\[
 H/B_s^c=CF_s^cR\in\Int(A).
\]
The same argument applied to $J$ proves the second assertion.
\end{proof}

\begin{proof}[Proof of Theorem~\ref{thm:main}]
The units of $\Int(A)$ are exactly $\F_q^\times$. Indeed, a polynomial unit and its inverse must both have $X$-degree zero, and the constant elements of $\Int(A)$ are exactly $A$.

Suppose $HJ=P$ with $H,J\in\Int(A)$. We induct on the largest index $s$ such that $a_s>0$, omitting trailing zero exponents. If all exponents are zero, $HJ=1$, and the assertion follows from the description of the units. If $s=0$, unique factorization in $K[X]$ gives
\[
 H=CX^b,\qquad J=C^{-1}X^{a_0-b},\qquad0\leq b\leq a_0.
\]
Evaluation at $X=1$ shows $C,C^{-1}\in A$, so $C\in\F_q^\times$.

Now let $s\geq1$. Apply Lemma~\ref{lem:extract}. The two resulting integer-valued quotients have product
\[
 \frac{H}{B_s^c}\frac{J}{B_s^{a_s-c}}
 =\prod_{i=0}^{s-1}B_i^{a_i}.
\]
The induction hypothesis gives
\[
 H/B_s^c=u\prod_{i=0}^{s-1}B_i^{b_i},
 \qquad u\in\F_q^\times,\quad0\leq b_i\leq a_i.
\]
Taking $b_s=c$ proves~\eqref{eq:divisor}. Conversely, a polynomial of that form has an integer-valued complementary product, so it divides $P$.

For uniqueness, compare the multiplicity of any root in $V_s\setminus V_{s-1}$ in two expressions of the form~\eqref{eq:divisor}. This determines $b_s$. Cancel this power of $B_s$ and descend. At the final level, the multiplicity of zero determines $b_0$. Equality of the remaining constants then determines $u$.
\end{proof}

\begin{proof}[Proof of Corollary~\ref{cor:question}]
The irreducibility of each $B_i$ was proved in~\cite{TW}; it also follows by applying Theorem~\ref{thm:main} to $P=B_i$. If an irreducible polynomial divides a product $P$ in that theorem, its expression~\eqref{eq:divisor} must contain exactly one basic factor, counted with multiplicity. Hence every irreducible divisor of $P$ is associate to some $B_i$. Uniqueness of the exponent vector proves the assertion after taking $a_i=m\alpha_i$.
\end{proof}

\section{Consequences and examples}

\begin{corollary}\label{cor:monoid}
The submonoid
\[
 \mathcal C=\left\{u\prod_{i\geq0}B_i^{a_i}:
 u\in\F_q^\times,\ a_i\in\Nzero,
 \text{ finitely many }a_i\neq0\right\}
\]
of $\Int(A)\setminus\{0\}$ is divisor-closed. Its quotient by its units is the free commutative monoid on the classes of $B_0,B_1,\ldots$.
\end{corollary}

\begin{proof}
Divisor closure is exactly Theorem~\ref{thm:main}. The uniqueness statement in that theorem shows that the exponent-vector map identifies $\mathcal C/\F_q^\times$ with the direct sum $\bigoplus_{i\geq0}\Nzero$.
\end{proof}

In particular, if $k=\sum_i\alpha_iq^i\geq1$, then the divisor classes of $\beta_k^m$ form a product of finite chains:
\[
 \{[H]:H\mid\beta_k^m\}
 \cong\prod_{i:\alpha_i>0}\{0,1,\ldots,m\alpha_i\}.
\]
This is also an isomorphism of partially ordered sets, with divisibility on the left and coordinatewise order on the right. Thus
\[
 \#\{[H]:H\mid\beta_k^m\}=\prod_i(m\alpha_i+1),
 \qquad
 \mathsf L(\beta_k^m)=\left\{m\sum_i\alpha_i\right\},
\]
where $\mathsf L$ denotes the set of factorization lengths. The uniqueness of factorization is stronger than the singleton assertion for lengths.

\begin{example}
For every $s>r\geq0$, the index $k=q^s+q^r$ gives
\[
 \beta_k=B_sB_r.
\]
Its $m$th power has exactly $(m+1)^2$ divisor classes, represented by $B_s^uB_r^v$ with $0\leq u,v\leq m$. Its unique factorization has length $2m$. These mixed-index families have unbounded degree and are not covered by the prime-power assertion alone.
\end{example}

\begin{remark}
The analogue of Corollary~\ref{cor:monoid} for the ordinary binomial polynomials fails. Put $b_n(X)=\binom Xn$. Then
\[
 b_2(X)b_3(X)=X\,G(X),\qquad
 G(X)=\frac{X(X-1)^2(X-2)}{12}\in\Int(\mathbb Z).
\]
The numerator of $G(x)$ is divisible by $3$, as it contains three consecutive integers. It is also divisible by $4$: for odd $x$ the square $(x-1)^2$ suffices, while for even $x$ the product $x(x-2)$ suffices. However, $G$ is not a product of ordinary binomial polynomials, up to a unit. Any such product with roots contained in $\{0,1,2\}$ has weakly decreasing root multiplicities in the order $0,1,2$, whereas those of $G$ are $(1,2,1)$. Thus products of the ordinary binomial polynomials need not form a divisor-closed monoid, even though each $b_n$ is absolutely irreducible~\cite{RW}.
\end{remark}

\begin{thebibliography}{9}

\bibitem{RW}
R. Rissner and D. Windisch,
\emph{Absolute irreducibility of the binomial polynomials},
Journal of Algebra \textbf{578} (2021), 92--114.
\href{https://doi.org/10.1016/j.jalgebra.2021.03.007}{doi:10.1016/j.jalgebra.2021.03.007}.

\bibitem{TW}
R. Tichy and D. Windisch,
\emph{Irreducibility properties of Carlitz' binomial coefficients for algebraic function fields},
Finite Fields and Their Applications \textbf{96} (2024), article 102413.
\href{https://doi.org/10.1016/j.ffa.2024.102413}{doi:10.1016/j.ffa.2024.102413}.
Preprint: \href{https://arxiv.org/abs/2310.02061}{arXiv:2310.02061}.

\bibitem{Wagner}
C. G. Wagner,
\emph{On the factorization of some polynomial analogues of binomial coefficients},
Archiv der Mathematik \textbf{24} (1973), 50--52.
\href{https://web.math.utk.edu/~wagner/papers/analogues.pdf}{Author's copy}.

\end{thebibliography}
\end{document}
